\documentclass[11pt]{article}
% Shared preamble for the Cupcast papers.
\usepackage[a4paper,margin=2.8cm]{geometry}
\usepackage{amsmath,amssymb,mathtools}
\usepackage{booktabs}
\usepackage{graphicx}
\usepackage{microtype}
% Classic BibTeX via natbib: tectonic runs bibtex natively, while its bundled
% biblatex is version-locked against system biber. Each paper folder symlinks
% references.bib from shared/.
\usepackage[numbers,sort&compress]{natbib}
\usepackage{xcolor}
\usepackage[colorlinks=true,allcolors=blue!45!black]{hyperref}
\graphicspath{{../figures/}}

% Shared notation
\newcommand{\E}{\mathbb{E}}
\newcommand{\Prob}{\mathbb{P}}
\newcommand{\att}{\alpha}
\newcommand{\dfn}{\beta}


\title{Cupcast~II-V: Held-Out Validation and a Comparison\\
Against the Cupcast~I Baseline on the 2026 World Cup}
\author{Carlos Garcia}
\date{June 2026}

\begin{document}
\maketitle

\begin{abstract}
We evaluate the dynamic state-space model of Paper~II against the static
Cupcast~I baseline on the matches of the 2026 World Cup that have actually been
played, both models trained strictly before kickoff. On the $60$ completed group
matches the new model improves out-of-sample log-loss from $0.869$ to $0.856$
and is better on Brier score and the ranked probability score as well. A
pre/post ablation shows the bounded context layer (travel, host, altitude)
improves all three scores further; an honest counter-result is that
cross-validated hyperparameter tuning, while it lowers the internationals
cross-validation loss, \emph{degrades} World Cup performance---the model's
default priors generalise better than tuned ones. We close with the headline
forecast.
\end{abstract}

\section{Protocol}
Both models are fitted to international results before the tournament cutoff of
11~June~2026; the $60$ subsequently completed group matches are a clean held-out
test set. We report three strictly proper scoring rules \cite{gneitingraftery2007}
over the three-way (home/draw/away) outcome: logarithmic loss, the multiclass
Brier score, and the ranked probability score (RPS). Lower is better; the uniform
forecast $(\tfrac13,\tfrac13,\tfrac13)$ scores $\log$-loss $\ln 3 = 1.099$.
The Cupcast~I baseline is its frozen pipeline, queried through a name-mapping
layer; the comparison is restricted to matches both models can score.

\section{The rebuild improves the metric}
Table~\ref{tab:headline} reports the held-out scores. The dynamic model beats the
static baseline on every metric, and both beat the uniform forecast decisively.

\begin{table}[h]\centering
\begin{tabular}{lccc}
\toprule
Forecaster & Log-loss & Brier & RPS \\
\midrule
Cupcast~II (dynamic + clubform) & \textbf{0.856} & \textbf{0.515} & \textbf{0.157} \\
Cupcast~I (static baseline)     & 0.869 & 0.523 & 0.160 \\
Uniform                         & 1.099 & 0.667 & 0.233 \\
\bottomrule
\end{tabular}
\caption{Held-out scoring on the 60 played 2026 World Cup group matches. Lower
is better.}\label{tab:headline}
\end{table}

The improvement is modest ($\sim1.4\%$ in log-loss) but consistent across all
three rules. It is corroborated out of sample on $772$ historical international
matches under rolling-origin cross-validation, where the dynamic model scores
log-loss $0.910$ (RPS $0.175$) against $1.099$ for uniform, on par with a
static weighted Dixon--Coles baseline. The most visible behavioural change is
that the Cupcast~I overrating of recent competitive form is corrected: the title
probability the static model assigned to its favourite falls from an inflated
$16.3\%$ to a market-plausible $10.2\%$.

\section{The context layer earns its place}
We ablate the bounded context covariates of Paper~II by scoring the same
held-out matches with and without them (host advantage routed entirely through
the context layer so the comparison is clean). Table~\ref{tab:ablation} shows
that travel, host, and altitude together improve all three scores.

\begin{table}[h]\centering
\begin{tabular}{lccc}
\toprule
 & Log-loss & Brier & RPS \\
\midrule
Model only            & 0.865 & 0.521 & 0.160 \\
Model + context       & \textbf{0.860} & \textbf{0.518} & \textbf{0.158} \\
\bottomrule
\end{tabular}
\caption{Pre/post-context ablation on the 60 played group matches.}\label{tab:ablation}
\end{table}

The signals are concrete: Mexico plays all three group games at home at altitude
(Mexico City $\approx2240$\,m, Guadalajara $\approx1560$\,m), and Iran faces a
$1{,}553$\,km travel leg between matches. Each effect is hard-capped and
provenance-logged, so the gain is earned without letting any single signal
dominate.

\section{An honest counter-result: tuning overfits the wrong regime}
We tuned the model's prior hyperparameters by Optuna \cite{akiba2019optuna}
under rolling-origin cross-validation on past internationals. The search found
parameters with a \emph{better} cross-validation loss ($0.842$ versus the
defaults), but those parameters \emph{worsened} the World Cup hold-out (log-loss
$0.865$, up from the untuned $0.856$). The internationals folds---qualifiers and
friendlies---are a different regime from the tournament (neutral venues,
fresh squads, knockout intensity), and tuning to them does not transfer.
The practical conclusion is to keep the sensible default priors, which are robust
and remain the best configuration on the actual test.

\section{The forecast}
The production configuration (default priors, model host advantage, and the
context travel/altitude covariates) yields the championship probabilities of
Table~\ref{tab:forecast} from $50{,}000$ simulated tournaments.

\begin{table}[h]\centering
\begin{tabular}{lc@{\qquad}lc}
\toprule
Team & Champion & Team & Champion \\
\midrule
Brazil    & 20.4\% & Netherlands & 5.5\% \\
Spain     & 13.2\% & Belgium     & 4.7\% \\
England   & 11.9\% & Colombia    & 4.2\% \\
Argentina & 10.2\% & Germany     & 2.6\% \\
Portugal  & 8.3\%  & Switzerland & 1.8\% \\
France    & 5.8\%  & Uruguay     & 1.5\% \\
\bottomrule
\end{tabular}
\caption{Cupcast~II championship probabilities for the 2026 World Cup
(top teams), from 50{,}000 Monte Carlo tournaments with the held-out model.}\label{tab:forecast}
\end{table}

\section{Discussion}
Two of the redesign's premises are validated on the real tournament: a
lineup-aware, dynamically-evolving model beats the static results-only baseline,
and bounded contextual signals add genuine predictive value. A third, more
sobering, lesson is that aggressive hyperparameter tuning can overfit a proxy
distribution; the disciplined defaults win. Remaining work---weather and injury
covariates, and a learned (rather than shadow) morale signal subject to the same
ablation gate---can only be admitted to the production forecast if it clears the
same held-out bar these results were held to.

\bibliographystyle{abbrvnat}
\bibliography{references}

\end{document}
